\section{Approach}\label{sec:approach}

Kairo constructs a feature matrix for each perspective, scales and optionally reduces it, assigns states by clustering, and compares state distributions over time. \autoref{fig:approach-overview} shows how the stages connect.

\subsection{Representing Process Behavior}\label{sec:approach-features}

The unit of observation depends on the perspective. Event-level observations retain position within a case. Window-level observations summarize conditions shared across cases over a calendar interval. \Cref{def:feature-representation} defines these observation sequences and the feature rows attached to them.

\begin{Definition}[Observation sequences and feature maps]\label{def:feature-representation}
Write $p$ for the selected perspective, with $p\in\{\mathrm{intra},\mathrm{res},\mathrm{inter}\}$.
\begin{itemize}[leftmargin=*,itemsep=.2em,topsep=.3em]
  \item \textbf{Intra-case observations.} The observation sequence is $\sigma_{\mathrm{intra}}=\sigma(E)$, with one observation per event and time anchor $t_{\mathrm{intra}}(e)=t(e)$. Its feature history comes from the trace of its own case.
  \item \textbf{Resource and inter-case observations.} Choose a width $\delta_p$ and use the covering grid segment $\sigma_p=\mathcal W_{\delta_p}[\sigma(E)]$ from \Cref{def:calendar-grid}. Each interval $W$ in this sequence gives one feature window, with anchor $t_p(W)=b(W)$. These sequences retain intervals without events.
  \item \textbf{Feature map.} Let $O_p=\operatorname{set}(\sigma_p)$ be the set of observations. The map $x_p\colon O_p\to\mathbb R^{d_p}$ assigns each observation a vector containing its $d_p$ selected feature coordinates.
  \item \textbf{Feature values.} For a proposition $B$, the indicator $\mathbf{1}[B]$ is one if $B$ holds and zero otherwise. Ratios with zero denominator, means of empty samples, and sample standard deviations based on fewer than two values are set to zero. Elapsed gaps and event durations are expressed in minutes.
\end{itemize}
\end{Definition}

The time anchor links each row to the sequence used for drift scoring in \Cref{sec:approach-signal}. The following definitions specify which events contribute to the feature coordinates.

\subsubsection{Intra-Case Perspective}

\Cref{def:intra-features} describes the current activity, earlier activities, and how far the case has progressed. Each event is evaluated in the order of its own trace.

\begin{Definition}[Intra-case features]\label{def:intra-features}
For $c_0\in C$, use the event subset $E_{c_0}$ and its ordered trace $\tau_{c_0}$ from \Cref{def:ordered-log}. For $e=e_{c_0,j}$, the feature history is the position window $\pi(e)=\tau_{c_0}[1:j]$ from \Cref{def:case-history}. For $a^\star\in A$, let $N_{a^\star}(e)=|\{v\in\operatorname{set}(\pi(e)):a(v)=a^\star\}|$.
\begin{itemize}[leftmargin=*,itemsep=.2em,topsep=.3em]
  \item \textbf{Current activity.} The coordinate for $a^\star$ is $\mathbf{1}[a(e)=a^\star]$.
  \item \textbf{Past activities.} For a selected integer lag limit $h\geq0$, use one indicator block for each $\ell\in\{1,\ldots,h\}$. The coordinate for $a^\star$ in block $\ell$ is $\mathbf{1}[a(e_{c_0,j-\ell})=a^\star]$ if $j>\ell$, and zero otherwise.
  \item \textbf{Prefix presence.} The indicator for $a^\star$ is $\mathbf{1}[N_{a^\star}(e)>0]$.
  \item \textbf{Prefix frequency.} The share of $a^\star$ is $N_{a^\star}(e)/j$.
  \item \textbf{Directly-follows counts.} For each pair $(a^\star,a')$ observed in the log, the number of its occurrences within the prefix is
  \[
    \bigl|\{v\in\operatorname{set}(\pi(e)):v^-\text{ exists},\
                 (a(v^-),a(v))=(a^\star,a')\}\bigr|.
  \]
  \item \textbf{Case progress.} The normalized position is $j/m_{c_0}$.
\end{itemize}
The selected groups of these coordinates form $x_{\mathrm{intra}}(e)$.
\end{Definition}

The prefix ends at the current event. The optional progress coordinate uses the complete recorded trace length, so it provides a retrospective measure of position.

\subsubsection{Resource Perspective}

\Cref{def:resource-features} groups resource activity into the feature windows from \Cref{def:feature-representation}. It assigns a handover to the interval containing the receiving event, which allows transfers across window boundaries to contribute.

\begin{Definition}[Resource features]\label{def:resource-features}
Let $E^{\mathrm{res}}=\{e\in E:r(e)\ne\bot\}$ be the events with recorded resources. For $W\in O_{\mathrm{res}}$, apply \Cref{def:calendar-grid} to $\sigma(E^{\mathrm{res}})$ and set $E_W^{\mathrm{res}}=\operatorname{set}(\sigma(E^{\mathrm{res}})_W)$. For $\rho\in R$, let $E_{W,\rho}=\{e\in E_W^{\mathrm{res}}:r(e)=\rho\}$. The receiving events of handovers are
\[
 H_W=\{e\in E_W^{\mathrm{res}}: e^-\text{ exists},\
                                  r(e^-)\notin\{\bot,r(e)\}\}.
\]
The predecessor is taken from the full case trace. For $\rho,\nu\in R$, let $H_{W,\rho,\nu}=\{e\in H_W:r(e^-)=\rho,\ r(e)=\nu\}$.
\begin{itemize}[leftmargin=*,itemsep=.2em,topsep=.3em]
  \item \textbf{Events per resource.} The count for $\rho$ is $|E_{W,\rho}|$.
  \item \textbf{Cases per resource.} The distinct case count is $|\{c(e):e\in E_{W,\rho}\}|$.
  \item \textbf{Mean event duration.} Average the available $\operatorname{dur}(e)$ over $e\in E_{W,\rho}$.
  \item \textbf{Mean incoming handover gap.} Average $g(e)$ over $e\in H_W\cap E_{W,\rho}$.
  \item \textbf{Activity-resource share.} For each observed pair $(a^\star,\rho)$, the share of that activity performed by $\rho$ is
  \[
     \frac{|\{e\in E_{W,\rho}:a(e)=a^\star\}|}
          {|\{e\in E_W^{\mathrm{res}}:a(e)=a^\star\}|}.
  \]
  \item \textbf{Handover share.} For each observed handover pair $(\rho,\nu)$, the share of transfers from $\rho$ received by $\nu$ is
  \[
     \frac{|H_{W,\rho,\nu}|}{\sum_{\kappa\in R}|H_{W,\rho,\kappa}|}.
  \]
\end{itemize}
The selected feature groups and resources determine $x_{\mathrm{res}}(W)$. Each retained handover-share coordinate uses two selected resources. The denominators in both shares use all recorded resources.
\end{Definition}

Activity shares can reveal a reassignment even when the total number of events stays stable. The full denominators preserve that interpretation when the analyst displays a subset of resources.

\subsubsection{Inter-Case Perspective}

\Cref{def:inter-features} uses the full event sequence to summarize activity and timing across cases. The distinction between case arrivals and completions helps separate changes in incoming demand from changes in how quickly work finishes.

\begin{Definition}[Inter-case features]\label{def:inter-features}
For $W\in O_{\mathrm{inter}}$, take the event set $E_W$ obtained by applying \Cref{def:calendar-grid} to the full sequence $\sigma(E)$. The cases observed in the window form $C_W=\{c(e):e\in E_W\}$.
\begin{itemize}[leftmargin=*,itemsep=.2em,topsep=.3em]
  \item \textbf{Active cases.} The count is $|C_W|$.
  \item \textbf{Arrivals and completions.} Count cases with first or last recorded timestamps in the interval, respectively. The counts are $|\{c_0\in C:s_{c_0}\in W\}|$ and $|\{c_0\in C:f_{c_0}\in W\}|$.
  \item \textbf{Events per active case.} The ratio is $|E_W|/|C_W|$.
  \item \textbf{Within-case gaps.} Take the mean and sample standard deviation of $(g(e):e\in E_W,\ e^-\text{ exists})$. A gap contributes to the window of its later event.
  \item \textbf{Stalled cases.} Given a stall threshold $\eta_{\mathrm{stall}}>0$, the number of cases stalled at the interval end is
  \[
     \bigl|\{c_0\in C:\exists\,e\in E_{c_0}\ \text{with}\ e^+\text{ defined and }
                 t(e)+\eta_{\mathrm{stall}}<d(W)<t(e^+)\}\bigr|.
  \]
  \item \textbf{Numerical attributes.} For each mapped numerical case attribute $\alpha$, take the mean and sample standard deviation of the available values $(v_\alpha(e):e\in E_W,\ v_\alpha(e)\ne\bot)$.
  \item \textbf{Categorical attributes.} For each observed value $z$ of a mapped categorical case attribute $\alpha$, its share is $|\{e\in E_W:v_\alpha(e)=z\}|/|E_W|$.
\end{itemize}
The selected groups form $x_{\mathrm{inter}}(W)$. Attribute summaries count each case once per event it contributes.
\end{Definition}

A stalled case may have no events in the current feature window. For categorical shares, the denominator includes events whose category is missing.

\subsubsection{Running Example}\label{sec:approach-example}

Consider a service desk observed over consecutive feature windows $W_{\mathrm a},\ldots,W_{\mathrm f}$. In each window $W_{\mathrm x}$, a focal case $\mathrm x$ and its companion case overlap in time. Each has a Receive event followed by a Resolve event, and both events fall within that window. The Frontline team receives every case and resolves every companion case after a gap of $g$. It also resolves the focal cases $\mathrm a$, $\mathrm b$, and $\mathrm c$ after the same gap. The Specialist team resolves the focal cases $\mathrm d$, $\mathrm e$, and $\mathrm f$ after a longer gap $G>g$.

\begin{Example}[Intra-case features for a case prefix]\label{ex:intra-case}
Compare the Receive and Resolve prefixes of cases $\mathrm c$ and $\mathrm d$. At Receive, the current-activity indicator for Receive is on and the prefix contains no directly-follows pair. At Resolve, Receive is the immediately preceding activity and the prefix contains the Receive-to-Resolve pair once. \autoref{tab:example-intra-case} shows the same feature values at the corresponding positions of both cases. \autoref{fig:example-intra-case} illustrates the prefixes used to build the rows.
\end{Example}

\begin{table*}[!t]
  \caption{Selected intra-case features for the Receive and Resolve events of cases $\mathrm c$ and $\mathrm d$. On and off denote true and false for an indicator.}
  \label{tab:example-intra-case}
  \centering
  \small
  \begin{tabularx}{\linewidth}{@{}ll*{3}{>{\centering\arraybackslash}X}@{}}
    \toprule
    \textbf{Case} & \textbf{Current event} & \textbf{Receive now} &
      \textbf{Receive immediately before} & \textbf{Receive-to-Resolve count} \\
    \midrule
    $\mathrm c$ & Receive & On & Off & Zero \\
    $\mathrm c$ & Resolve & Off & On & One \\
    $\mathrm d$ & Receive & On & Off & Zero \\
    $\mathrm d$ & Resolve & Off & On & One \\
    \bottomrule
  \end{tabularx}
\end{table*}

\begin{figure*}[!t]
  \centering
  \begin{tikzpicture}[
  kairo/base, x=1cm, y=1cm,
  activity/.style={
    rounded corners=2pt, minimum width=1.58cm, minimum height=0.37cm,
    draw=KairoLine, fill=white, inner sep=2pt
  },
  current/.style={activity, draw=KairoBlue, fill=KairoBlueTint, text=KairoBlue},
  follows/.style={-{Stealth[length=4pt]}, draw=KairoMuted, line width=0.45pt}
]
  \path[use as bounding box] (0,0) rectangle (13.08,-1.56);
  \path[fill=KairoPaper, rounded corners=3pt] (1.69,-1.55) rectangle (7.23,-0.40);
  \path[fill=KairoPaper, rounded corners=3pt] (7.49,-1.55) rectangle (13.03,-0.40);
  \node[anchor=west, font=\footnotesize\bfseries] at (0,-0.16) {Case};
  \node[font=\footnotesize\bfseries] at (4.46,-0.16) {Prefix at Receive};
  \node[font=\footnotesize\bfseries] at (10.26,-0.16) {Prefix at Resolve};
  \foreach \caseid/\rowy in {c/-0.72,d/-1.23} {
    \node[anchor=west] at (0,\rowy) {$\mathrm{\caseid}$};
    \node[current] at (4.46,\rowy) {Receive};
    \node[activity, text=KairoMuted] (receive-\caseid) at (9.29,\rowy) {Receive};
    \node[current] (resolve-\caseid) at (11.23,\rowy) {Resolve};
    \draw[follows] (receive-\caseid.east) -- (resolve-\caseid.west);
  }
\end{tikzpicture}

  \caption{Cases $\mathrm c$ and $\mathrm d$ have matching activity prefixes at each event. The shaded activity is the current event in each prefix.}
  \label{fig:example-intra-case}
\end{figure*}

\begin{Example}[Resource features for a handover]\label{ex:resource}
Between $W_{\mathrm c}$ and $W_{\mathrm d}$, responsibility for resolving the focal case moves from Frontline to Specialist. Frontline still resolves the companion case, so Specialist handles half of the Resolve events in $W_{\mathrm d}$. The only handover in that window goes from Frontline to Specialist, with a gap of $G$. \autoref{tab:example-resource} shows how the assignment and handover features change. \autoref{fig:example-resource} shows who performs each focal event.
\end{Example}

\begin{table*}[!t]
  \caption{Selected resource features in consecutive feature windows. The handover share and mean incoming gap are zero in $W_{\mathrm c}$, which contains no handover.}
  \label{tab:example-resource}
  \centering
  \small
  \begin{tabularx}{\linewidth}{@{}X*{2}{>{\centering\arraybackslash}p{0.17\linewidth}}@{}}
    \toprule
    \textbf{Feature} & $\boldsymbol{W_{\mathrm c}}$ & $\boldsymbol{W_{\mathrm d}}$ \\
    \midrule
    Events by Frontline & Four & Three \\
    Events by Specialist & Zero & One \\
    Share of Resolve performed by Specialist & Zero & Half \\
    Share of Frontline's handovers received by Specialist & Zero & Full \\
    Mean incoming handover gap for Specialist & Zero & $G$ \\
    \bottomrule
  \end{tabularx}
\end{table*}

\begin{figure*}[!t]
  \centering
  \begin{tikzpicture}[
  kairo/base, x=1cm, y=1cm,
  receive/.style={circle, minimum size=4.3pt, inner sep=0pt, line width=0.75pt, fill=white},
  resolve/.style={rectangle, minimum size=4.1pt, inner sep=0pt},
  same team/.style={draw=KairoBlue, line width=0.85pt, line cap=round},
  handover/.style={draw=KairoAmber, line width=0.85pt, dash pattern=on 3pt off 1.5pt,
    -{Stealth[length=4pt]}, shorten >=3pt}
]
  \path[use as bounding box] (0,0) rectangle (13.08,-2.75);
  % Use the same local time positions as the inter-case diagram.
  \foreach \windowid/\leftx/\windowtitle/\focalend/\focaly in {
    c/2.10/Earlier/2.84/-1.31,
    d/7.81/Later/4.49/-1.97
  } {
    \begin{scope}[shift={(\leftx,0)}]
      \node[font=\footnotesize\bfseries] at (2.635,-0.17)
        {\windowtitle\ window $W_{\mathrm{\windowid}}$};
      \path[fill=KairoBlueTint, rounded corners=2pt] (0,-1.58) rectangle (5.27,-0.50);
      \path[fill=KairoAmberTint, rounded corners=2pt] (0,-2.26) rectangle (5.27,-1.64);
      \coordinate (companion-start-\windowid) at (1.46,-0.81);
      \coordinate (companion-end-\windowid) at (3.56,-0.81);
      \coordinate (focal-start-\windowid) at (0.74,-1.31);
      \coordinate (focal-end-\windowid) at (\focalend,\focaly);
    \end{scope}
  }
  \node[anchor=east, font=\footnotesize\bfseries] at (1.96,-0.49) {Frontline};
  \node[anchor=east, text=KairoMuted] at (1.96,-0.81) {Companion};
  \node[anchor=east, text=KairoMuted] at (1.96,-1.31) {Focal case};
  \node[anchor=east, font=\footnotesize\bfseries] at (1.96,-1.97) {Specialist};
  \foreach \windowid in {c,d} {
    \draw[same team] (companion-start-\windowid) -- (companion-end-\windowid);
    \node[receive, draw=KairoBlue] at (companion-start-\windowid) {};
    \node[resolve, fill=KairoBlue] at (companion-end-\windowid) {};
  }
  \draw[same team] (focal-start-c) -- (focal-end-c);
  \draw[handover] (focal-start-d) -- (focal-end-d);
  \node[anchor=south, text=KairoAmber, inner sep=1pt] at (10.69,-1.59) {Handover};
  \node[receive, draw=KairoBlue] at (focal-start-c) {};
  \node[receive, draw=KairoBlue] at (focal-start-d) {};
  \node[resolve, fill=KairoBlue] at (focal-end-c) {};
  \node[resolve, fill=KairoAmber] at (focal-end-d) {};
  \node[receive, draw=KairoMuted] at (4.85,-2.60) {};
  \node[anchor=west, text=KairoMuted] at (4.98,-2.60) {Receive};
  \node[resolve, fill=KairoMuted] at (6.56,-2.60) {};
  \node[anchor=west, text=KairoMuted] at (6.69,-2.60) {Resolve};
  \draw[-{Stealth[length=4pt]}, draw=KairoMuted, line width=0.5pt]
    (8.19,-2.60) -- (8.74,-2.60);
  \node[anchor=west, text=KairoMuted] at (8.84,-2.60) {Time within window};
\end{tikzpicture}

  \caption{The focal case passes from Frontline to Specialist in $W_{\mathrm d}$. The handover is assigned to the window containing the Specialist event.}
  \label{fig:example-resource}
\end{figure*}

\begin{Example}[Inter-case features for a longer gap]\label{ex:inter-case}
Both $W_{\mathrm c}$ and $W_{\mathrm d}$ have the same numbers of arrivals, completions, and active cases. Each Resolve event contributes its gap from Receive to the mean for its window. The gaps are both $g$ in $W_{\mathrm c}$. In $W_{\mathrm d}$, the focal case contributes $G$ and the companion still contributes $g$, so the mean rises. \autoref{tab:example-inter-case} shows this change in timing alongside the matching counts. \autoref{fig:example-inter-case} shows which gaps enter each window's mean.
\end{Example}

\begin{table*}[!t]
  \caption{Selected inter-case features in consecutive feature windows. The later window has the same counts and a longer mean gap.}
  \label{tab:example-inter-case}
  \centering
  \small
  \begin{tabularx}{\linewidth}{@{}X*{2}{>{\centering\arraybackslash}p{0.17\linewidth}}@{}}
    \toprule
    \textbf{Feature} & $\boldsymbol{W_{\mathrm c}}$ & $\boldsymbol{W_{\mathrm d}}$ \\
    \midrule
    Case arrivals & Two & Two \\
    Case completions & Two & Two \\
    Active cases & Two & Two \\
    Events per active case & Two & Two \\
    Mean within-case event gap & $g$ & $\operatorname{mean}(g,G)$ \\
    \bottomrule
  \end{tabularx}
\end{table*}

\begin{figure*}[!t]
  \centering
  \begin{tikzpicture}[
  kairo/base, x=1cm, y=1cm,
  receive/.style={circle, minimum size=4.3pt, inner sep=0pt, draw=KairoTeal,
    line width=0.75pt, fill=white},
  resolve/.style={rectangle, minimum size=4.1pt, inner sep=0pt, fill=KairoTeal},
  gap/.style={draw=KairoTeal, line width=0.9pt, line cap=round, text=KairoTeal}
]
  \path[use as bounding box] (0,0) rectangle (13.08,-2.20);
  % Use the same local time positions as the resource diagram.
  \foreach \windowid/\leftx/\windowtitle/\focalend/\focalgap/\windowmean in {
    c/2.10/Earlier/2.84/g/g,
    d/7.81/Later/4.49/G/{\operatorname{mean}(g,G)}
  } {
    \begin{scope}[shift={(\leftx,0)}]
      \node[font=\footnotesize\bfseries] at (2.635,-0.17)
        {\windowtitle\ window $W_{\mathrm{\windowid}}$};
      \path[fill=KairoPaper, rounded corners=2pt] (0,-1.66) rectangle (5.27,-0.50);
      \draw[gap] (1.46,-0.91) -- node[midway, above=1pt, inner sep=1pt] {$g$} (3.56,-0.91);
      \draw[gap] (0.74,-1.43) -- node[midway, above=1pt, inner sep=1pt] {$\focalgap$}
        (\focalend,-1.43);
      \node[receive] at (1.46,-0.91) {};
      \node[resolve] at (3.56,-0.91) {};
      \node[receive] at (0.74,-1.43) {};
      \node[resolve] at (\focalend,-1.43) {};
      \path[fill=KairoTealTint, rounded corners=2pt] (0,-2.15) rectangle (5.27,-1.77);
      \node[text=KairoTeal] at (2.635,-1.96) {$\windowmean$};
    \end{scope}
  }
  \node[anchor=east] at (1.96,-0.91) {Companion};
  \node[anchor=east] at (1.96,-1.43) {Focal case};
  \node[anchor=east, font=\footnotesize\bfseries] at (1.96,-1.96) {Window mean};
\end{tikzpicture}

  \caption{The focal case's longer gap raises the window mean in $W_{\mathrm d}$. Open circles mark arrivals and filled squares mark completions.}
  \label{fig:example-inter-case}
\end{figure*}


\subsection{Scaling and Reducing the Vectors}\label{sec:approach-reduction}

Feature columns can measure counts, proportions, or durations. A column with a large numerical range can dominate clustering distances. \Cref{def:cluster-representation} specifies how Kairo rescales columns and optionally projects observations to fewer dimensions. Case identifiers and timestamps remain attached to each row as metadata.

\begin{Definition}[Scaling and projection]\label{def:cluster-representation}
Fix a perspective $p$. The representation used for clustering and its projection diagnostics are defined as follows.
\begin{itemize}[leftmargin=*,itemsep=.2em,topsep=.3em]
  \item \textbf{Column scaling.} Write $x_{pi}(o)$ for coordinate $i$ of $x_p(o)$. A scaled coordinate is
  \[
     \widetilde{x}_{pi}(o)=\frac{x_{pi}(o)-\mu_{pi}}{\sigma_{pi}},
  \]
  where $\mu_{pi}$ and $\sigma_{pi}$ are the column mean and sample standard deviation for standard score scaling, or the column minimum and range for min-max scaling. A constant column selected for scaling maps to zero. Excluded columns keep their original values.
  \item \textbf{Principal component projection.} Let $\overline{x}_p$ be the mean of the scaled vectors and let $U_\ell$ contain the leading $\ell$ orthonormal covariance eigenvectors in decreasing order of eigenvalue. With principal component analysis (PCA), $z_p(o)=U_\ell^\top(\widetilde{x}_p(o)-\overline{x}_p)$. Without projection, $z_p(o)=\widetilde{x}_p(o)$.
  \item \textbf{Projection diagnostics.} Let $\xi_1,\ldots,\xi_{d_p}$ be the covariance eigenvalues in the same order. For positive total variance, the retained fraction is $\sum_{i=1}^{\ell}\xi_i/\sum_{i=1}^{d_p}\xi_i$. The mean squared reconstruction error is
  \[
     \frac{1}{|O_p|}\sum_{o\in O_p}
       \bigl\|\widetilde{x}_p(o)-\overline{x}_p-U_\ell z_p(o)\bigr\|^2.
  \]
\end{itemize}
\end{Definition}

The activity indicators can retain their original scale. PCA can compress related coordinates, such as activity presence and frequency~\cite{hotelling1933components}. Retained variance and reconstruction error help the analyst choose the number of components. A discarded direction can still carry evidence of change.

\subsection{Learning States}\label{sec:approach-states}

Clustering groups similar observations into recurring patterns. \Cref{def:state-model} specifies how labels describe the groups and how their order gives a trajectory.

\begin{Definition}[State model and trajectories]\label{def:state-model}
Fix a perspective $p$ and the clustering representation of \Cref{def:cluster-representation}.
\begin{itemize}[leftmargin=*,itemsep=.2em,topsep=.3em]
  \item \textbf{Labels and states.} A fitted model assigns a label $s_p(o)$ to each $o\in O_p$ using its representation $z_p(o)$. The observed label set is $S_p=\{s_p(o):o\in O_p\}$. The state for $s\in S_p$ consists of $\{o\in O_p:s_p(o)=s\}$.
  \item \textbf{Trajectory.} Applying the labels along a case trace $\tau_{c_0}$ gives an intra-case trajectory. Applying them along $\sigma_p$ gives the trajectory of a window-based perspective.
  \item \textbf{State visit.} A visit is a maximal contiguous subsequence with a common label. It begins at the first observation's anchor and ends at the next visit's start. The final visit ends at the last event's timestamp for a case, or at the end of the final feature window for a window-based trajectory.
\end{itemize}
\end{Definition}

Kairo fits the model once per perspective for the analysis presented here. The resulting labels are comparable across periods, and each fit uses the complete log. \Cref{def:clustering-rules} gives the assignment rules.

\begin{Definition}[Clustering rules]\label{def:clustering-rules}
Let $z=z_p(o)$ be the representation of an observation $o\in O_p$.
\begin{itemize}[leftmargin=*,itemsep=.2em,topsep=.3em]
  \item \textbf{Distances.} For vectors $z$ and $v$, the Euclidean distance is $d_{\mathrm E}(z,v)=\|z-v\|$. For nonzero vectors, the cosine distance is $d_{\mathrm C}(z,v)=1-\langle z,v\rangle/(\|z\|\|v\|)$. The function $d_{\mathrm{cl}}$ denotes the distance selected for a method that permits a choice.
  \item \textbf{$k$-means.} The model assigns each observation to its closest center in Euclidean distance. Fitting alternates assignment with moving each occupied center to the mean of its assigned vectors~\cite{lloyd1982quantization}. For $k$ centers $\mu_1,\ldots,\mu_k$, the objective is to minimize
  \[
      \sum_{o\in O_p}\|z_p(o)-\mu_{s_p(o)}\|^2 .
  \]
  \item \textbf{Self-organizing map.} A self-organizing map (SOM) puts a prototype $\mu_g$ at each grid cell $g\in G$. Its label for $z$ is a winning cell $g^\ast$ that minimizes $d_{\mathrm{cl}}(z,\mu_g)$. In training, a prototype takes the step $\beta H(g,g^\ast)(z-\mu_g)$, where $\beta$ is the learning rate and $H$ weights cells by their proximity to the winner~\cite{kohonen1982maps}.
  \item \textbf{Density-based clustering.} Density-based spatial clustering of applications with noise (DBSCAN) uses a radius $\varepsilon$ and a minimum count $m$. A vector is a core if the sum of multiplicities of distinct vectors within distance $\varepsilon$, including itself, is at least $m$. Connected components of neighboring core vectors form clusters. A non-core vector in a core neighborhood joins a cluster of such a core. The remaining vectors receive a noise label~\cite{ester1996density}. Kairo maps each label back to the original observations with that vector.
\end{itemize}
\end{Definition}

The number of centers controls how finely $k$-means can split the observations. The SOM grid supports inspection of neighboring prototypes. For DBSCAN, the dashboard shows the ordered distances at which each observation's neighborhood reaches the chosen count, including the observation itself. A fit that finds a cluster can include noise in its state shares.

\autoref{fig:approach-clustering} uses a schematic resource-feature cloud to illustrate the assignment rules. The cloud has variations around the Frontline and Specialist resolution patterns from the running example, along with an isolated observation.

\begin{figure}[htbp]
  \centering
  \subfloat[$k$-means\label{fig:approach-clustering-kmeans}]{%
    \FigurePlaceholder[0.28]{TikZ sketch of the shared points, fitted centroids, and nearest-centroid assignments.}}
  \hfill
  \subfloat[Self-organizing map\label{fig:approach-clustering-som}]{%
    \FigurePlaceholder[0.28]{TikZ sketch of the shared points linked to their winning cells on a grid with neighboring prototypes.}}
  \hfill
  \subfloat[DBSCAN\label{fig:approach-clustering-dbscan}]{%
    \FigurePlaceholder[0.28]{TikZ sketch of core neighborhoods and border points. Mark the isolated observation as noise.}}
  \caption{State assignment on a shared schematic feature cloud. Each point represents a resource feature window. The panels show the nearest centroid for $k$-means, the winning cell for a self-organizing map, and density-connected groups and noise for DBSCAN.}
  \label{fig:approach-clustering}
\end{figure}


An analyst can interpret a state by examining the original feature values of its members. Plots of between-state distances use cluster centers for $k$-means, cell prototypes for a SOM, and means of distinct core vectors for DBSCAN. These representatives support visualization. The drift signal uses observation counts by state.

\paragraph{Inspecting the fitted states.}
\Cref{def:cluster-diagnostics} gives the measures used to check how the model organizes observations.

\begin{Definition}[State diagnostics]\label{def:cluster-diagnostics}
\leavevmode
\begin{itemize}[leftmargin=*,itemsep=.2em,topsep=.3em]
  \item \textbf{SOM quantization error.} Average $d_{\mathrm{cl}}(z_p(o),\mu_{s_p(o)})$ over $o\in O_p$.
  \item \textbf{SOM occupancy entropy.} Let $f_g$ denote the share of observations assigned to cell $g$ of grid $G$. The normalized entropy is $-\sum_{g:f_g>0}f_g\log f_g/\log|G|$ if $|G|>1$, and zero for a single-cell grid.
  \item \textbf{Silhouette.} For observation $o$, let $a_o$ be its mean distance to other members of its state and $b_o$ its least mean distance to the members of another state. The score averages $(b_o-a_o)/\max(a_o,b_o)$ over the evaluated observations, with a contribution of zero for a singleton state or a zero denominator~\cite{rousseeuw1987silhouettes}.
  \item \textbf{Calinski-Harabasz score.} Let $n$ be the number of evaluated observations, $k$ the number of their occupied states, and $n_s$ the number in state $s$. Let $\bar z_s$ be the mean vector of state $s$ and $\bar z$ the overall mean. When $2\leq k<n$ and within-state variation is positive, the score is
  \[
     \frac{\sum_s n_s\|\bar z_s-\bar z\|^2/(k-1)}
          {\sum_o\|z_p(o)-\bar z_{s_p(o)}\|^2/(n-k)}.
  \]
  The numerator measures separation of the state means and the denominator measures variation within states~\cite{calinski1974dendrite}.
\end{itemize}
The silhouette and Calinski-Harabasz calculations exclude DBSCAN noise.
\end{Definition}

These measures describe the fitted representation. \Cref{sec:assessment} evaluates how well the resulting drift signals locate annotated changes.

\subsection{Measuring Changes in State Frequencies}\label{sec:approach-signal}

The detector builds a time series from \emph{signal windows}. Each signal window collects the labels of observations whose time anchors fall inside it. \Cref{def:state-distributions} gives the grouping rule and the associated proportions. Short signal windows improve time resolution but may contain too few observations to estimate proportions reliably. \autoref{fig:approach-signal-windows} illustrates how resource feature windows contribute labels to the signal.

\begin{Definition}[Signal windows and state distributions]\label{def:state-distributions}
For perspective $p$, use the observation sequence $\sigma_p$ and time anchors $t_p$ from \Cref{def:feature-representation}.
\begin{itemize}[leftmargin=*,itemsep=.2em,topsep=.3em]
  \item \textbf{Signal-window membership.} Apply \Cref{def:calendar-grid} to $\sigma_p$ with anchor map $t_p$ and width $\Delta>0$. For a window-based perspective require $\Delta\geq\delta_p$. Write $V_1,\ldots,V_m$ for the intervals with a nonempty observation subsequence, in order of start time.
  \item \textbf{Observed state shares.} Let $O_{p,j}=\operatorname{set}((\sigma_p)_{V_j})$ and $N_j=|O_{p,j}|$. Using the labels from \Cref{def:state-model}, the distribution on $S_p$ in interval $V_j$ is
  \[
     P_j(s)=\frac{|\{o\in O_{p,j}:s_p(o)=s\}|}{N_j},
     \qquad s\in S_p.
  \]
\end{itemize}
\end{Definition}

A feature window with an empty event subsequence can still contribute a state label because it has a feature row. A signal interval with no observation anchors has no state distribution.

The reference determines what counts as recent or typical behavior. \Cref{def:drift-reference} describes a local comparison with the previous populated window, a smoothed comparison with several recent windows, and a fixed retrospective comparison with the whole log.

\begin{Definition}[Reference distribution]\label{def:drift-reference}
For a populated signal window $V_j$ from \Cref{def:state-distributions}, let $I_j$ be the reference index set.
\begin{itemize}[leftmargin=*,itemsep=.2em,topsep=.3em]
  \item \textbf{Previous window.} Use $I_j=\{j-1\}$ if $j>1$.
  \item \textbf{Recent windows.} For an integer lookback $h\geq1$, use $I_j=\{j-h,\ldots,j-1\}$ if $j>h$.
  \item \textbf{Whole log.} Use all populated indices, $I_j=\{1,\ldots,m\}$, including the current index.
\end{itemize}
When the index set is available, the reference is the unweighted mean $Q_j=|I_j|^{-1}\sum_{i\in I_j}P_i$.
\end{Definition}

\Cref{def:drift-score} compares an observed distribution with its reference. Each comparison places the resulting score at the start of the window it describes.

\begin{Definition}[Divergence and drift score]\label{def:drift-score}
For distributions $P,Q$ on $S_p$, let $M=(P+Q)/2$ and let $\epsilon>0$ be a fixed smoothing constant.
\begin{itemize}[leftmargin=*,itemsep=.2em,topsep=.3em]
  \item \textbf{Smoothed Kullback-Leibler term.} Define
  \[
    K_\epsilon(P,Q)=
      \sum_{s\in S_p}(P(s)+\epsilon)\log\frac{P(s)+\epsilon}{Q(s)+\epsilon}.
  \]
  \item \textbf{Distribution comparisons.} With square roots taken componentwise, the available functions are
  \[
    \begin{aligned}
      D_{\mathrm{KL}}(P,Q)&=K_\epsilon(P,Q),\\
      D_{\mathrm{JS}}(P,Q)&=\bigl(K_\epsilon(P,M)+K_\epsilon(Q,M)\bigr)/2,\\
      D_{\mathrm{TV}}(P,Q)&=\tfrac12\sum_{s\in S_p}|P(s)-Q(s)|,\\
      D_{\mathrm H}(P,Q)&=\|\sqrt P-\sqrt Q\|/\sqrt{2}.
    \end{aligned}
  \]
  \item \textbf{Drift score.} For a chosen function $D$ and an available reference $Q_j$, the score at $b(V_j)$ is $\lambda_j=D(P_j,Q_j)$. The score is undefined when the reference is unavailable.
\end{itemize}
\end{Definition}

These comparisons are smoothed Kullback-Leibler divergence, Jensen-Shannon divergence, total variation, and Hellinger distance. The smoothing keeps the logarithmic terms finite when a state is absent in a window. Jensen-Shannon averages differences to a common midpoint~\cite{lin1991divergence}. Total variation sums absolute changes in state shares. Hellinger uses the Euclidean distance between the vectors of square-root shares, with normalization.

\begin{Example}[From state labels to a drift score]
\label{ex:drift-score}
Consider the resource perspective of the service-desk log in \Cref{ex:service-desk}. Suppose the fitted model assigns $W_{\mathrm a},W_{\mathrm b},W_{\mathrm c}$ to state $\mathsf R$, in which Frontline resolves both cases. It assigns $W_{\mathrm d},W_{\mathrm e},W_{\mathrm f}$ to state $\mathsf S$, in which Specialist resolves the focal case. Choose a signal width twice the feature width, with $b(W_{\mathrm a})$ on a signal-window boundary. The grouping rule in \Cref{def:state-distributions} yields
\[
  V_{\mathrm A}=W_{\mathrm a}\cup W_{\mathrm b},\qquad
  V_{\mathrm B}=W_{\mathrm c}\cup W_{\mathrm d},\qquad
  V_{\mathrm C}=W_{\mathrm e}\cup W_{\mathrm f},
\]
as depicted in \autoref{fig:approach-signal-windows}.

The state distributions $P_{\mathrm A}$, $P_{\mathrm B}$, and $P_{\mathrm C}$ count feature-window labels within the respective signal windows and divide by the number of labels. All labels in $V_{\mathrm A}$ are $\mathsf R$, the labels in $V_{\mathrm B}$ are evenly split, and all labels in $V_{\mathrm C}$ are $\mathsf S$ (\autoref{fig:approach-state-distributions}). The middle window mixes the earlier and later patterns because it spans the change.

To score $V_{\mathrm C}$, the previous-window reference is $Q_{\mathrm C}^{\mathrm{prev}}=P_{\mathrm B}$. The recent-window reference that uses both earlier signal windows is $Q_{\mathrm C}^{\mathrm{recent}}=\operatorname{mean}(P_{\mathrm A},P_{\mathrm B})$. Its share of $\mathsf S$ is smaller than in $P_{\mathrm B}$ because $V_{\mathrm A}$ contains only $\mathsf R$. The whole-log reference also includes $P_{\mathrm C}$. Its mean is $P_{\mathrm B}$ here because the exclusive $\mathsf R$ and $\mathsf S$ distributions balance each other (\autoref{fig:approach-reference-distributions}).

For these states, total variation reduces to the absolute difference in the share of $\mathsf S$. Write $\alpha=P_{\mathrm B}(\mathsf S)$ for the middle window's share of this state. This gives equal previous-window scores at $b(V_{\mathrm B})$ and $b(V_{\mathrm C})$.
\[
  D_{\mathrm{TV}}(P_{\mathrm B},P_{\mathrm A})
  =D_{\mathrm{TV}}(P_{\mathrm C},P_{\mathrm B})
  =\alpha.
\]
The recent-window reference gives $V_{\mathrm C}$ a higher score because its reference contains less $\mathsf S$. The whole-log mean gives $V_{\mathrm B}$ a zero score and the same positive score to both outer windows. \autoref{fig:approach-drift-scores} places each value at the start of the window it describes. Local scores are undefined until the required earlier signal windows are available.
\end{Example}

\begin{figure*}[!t]
  \centering
  \begin{tikzpicture}[
  kairo/base,
  x=0.93cm,
  row label/.style={anchor=east, align=right, inner sep=0pt, text=KairoMuted},
  window edge/.style={draw=KairoLine, line width=0.45pt},
  signal edge/.style={draw=KairoMuted, line width=0.55pt}
]
  \node[row label] at (-0.21,1.66) {Signal window};
  \node[row label] at (-0.21,1.04) {Feature window};
  \node[row label] at (-0.21,0.45) {State at start};

  \foreach \fromX/\toX/\windowName/\shade in
    {0/4/A/KairoPaper,4/8/B/KairoTealTint,8/12/C/KairoPaper} {
    \fill[fill=\shade] (\fromX,1.41) rectangle (\toX,1.91);
    \draw[signal edge] (\fromX,1.41) rectangle (\toX,1.91);
    \node at ({(\fromX+\toX)/2},1.66) {$V_{\mathrm{\windowName}}$};
  }

  \foreach \fromX/\windowName/\state/\shade/\ink in
    {0/a/R/KairoBlueTint/KairoBlue,2/b/R/KairoBlueTint/KairoBlue,
     4/c/R/KairoBlueTint/KairoBlue,6/d/S/KairoAmberTint/KairoAmber,
     8/e/S/KairoAmberTint/KairoAmber,10/f/S/KairoAmberTint/KairoAmber} {
    \fill[fill=\shade] (\fromX,0.80) rectangle ++(2,0.48);
    \draw[window edge] (\fromX,0.80) rectangle ++(2,0.48);
    \node at ({\fromX+1},1.04) {$W_{\mathrm{\windowName}}$};
    \draw[window edge] (\fromX,0.80) -- (\fromX,0.45);
    \fill[fill=\ink] (\fromX,0.45) circle[radius=1.8pt];
    \node[anchor=west, inner sep=2.5pt, text=\ink] at (\fromX,0.45)
      {$\mathsf{\state}$};
  }
  \foreach \boundary in {0,4,8,12} {
    \draw[signal edge] (\boundary,1.41) -- (\boundary,0.80);
  }
  \draw[KairoMuted, -{Latex[length=1.65mm]}, line width=0.45pt]
    (0,0.08) -- (12,0.08) node[above left, inner sep=1.5pt] {Time};
\end{tikzpicture}

  \caption{Signal windows in the service-desk example. In the middle signal window, Frontline resolves both cases in the earlier feature window and Specialist resolves the focal case in the later one.}
  \label{fig:approach-signal-windows}
\end{figure*}

\begin{figure*}[!t]
  \centering
  \begin{tikzpicture}[kairo/base, x=1cm, y=1cm]
  \foreach \fromX/\windowName/\barDescription in {
    0/A/{Frontline only},
    4.50/B/{Equal shares},
    9.00/C/{Specialist resolves the focal case}%
  } {
    \node[anchor=west, inner sep=0pt] at (\fromX,1.32) {$P_{\mathrm{\windowName}}$};
    \node[anchor=east, inner sep=0pt, text=KairoMuted] at ({\fromX+4.10},1.32)
      {Window $V_{\mathrm{\windowName}}$};
    \node[anchor=north, inner sep=0pt, text=KairoMuted] at ({\fromX+2.05},0.53)
      {\barDescription};
  }

  \foreach \fromX/\toX/\state/\shade/\ink in {
    0/4.10/R/KairoBlueTint/KairoBlue,
    4.50/6.55/R/KairoBlueTint/KairoBlue,
    6.55/8.60/S/KairoAmberTint/KairoAmber,
    9.00/13.10/S/KairoAmberTint/KairoAmber%
  } {
    \fill[fill=\shade] (\fromX,0.65) rectangle (\toX,1.10);
    \draw[draw=\ink, line width=0.5pt] (\fromX,0.65) rectangle (\toX,1.10);
    \node[text=\ink] at ({(\fromX+\toX)/2},0.875) {$\mathsf{\state}$};
  }
\end{tikzpicture}

  \caption{Observed state distributions in the signal windows. The first and last windows contain only $\mathsf R$ and $\mathsf S$, respectively. The middle window contains equal shares.}
  \label{fig:approach-state-distributions}
\end{figure*}

\begin{figure*}[!t]
  \centering
  \begin{tikzpicture}[kairo/base, x=1cm, y=1cm]
  \foreach \fromX/\barTitle/\barSource/\ink in {
    0/{Current window}/{$P_{\mathrm C}$}/KairoInk,
    3.52/{Previous window}/{$P_{\mathrm B}$}/KairoMuted,
    6.88/{Recent mean}/{$\operatorname{mean}(P_{\mathrm A},P_{\mathrm B})$}/KairoTeal,
    10.24/{Whole-log mean}/{Equals $P_{\mathrm B}$ here}/KairoMuted%
  } {
    \node[inner sep=0pt, text=\ink, font=\footnotesize\bfseries]
      at ({\fromX+1.435},1.32) {\barTitle};
    \node[anchor=north, inner sep=0pt, text=KairoMuted]
      at ({\fromX+1.435},0.53) {\barSource};
  }
  \draw[KairoLine, line width=0.5pt] (3.19,0.22) -- (3.19,1.45);

  \foreach \fromX/\toX/\state/\shade/\ink in {
    0/2.87/S/KairoAmberTint/KairoAmber,
    3.52/4.955/R/KairoBlueTint/KairoBlue,
    4.955/6.39/S/KairoAmberTint/KairoAmber,
    6.88/9.0325/R/KairoBlueTint/KairoBlue,
    9.0325/9.75/S/KairoAmberTint/KairoAmber,
    10.24/11.675/R/KairoBlueTint/KairoBlue,
    11.675/13.11/S/KairoAmberTint/KairoAmber%
  } {
    \fill[fill=\shade] (\fromX,0.65) rectangle (\toX,1.10);
    \draw[draw=\ink, line width=0.5pt] (\fromX,0.65) rectangle (\toX,1.10);
    \node[text=\ink] at ({(\fromX+\toX)/2},0.875) {$\mathsf{\state}$};
  }
\end{tikzpicture}

  \caption{Reference distributions for scoring $V_{\mathrm C}$. The recent-window reference has less $\mathsf S$ than the previous-window and whole-log references in this example.}
  \label{fig:approach-reference-distributions}
\end{figure*}

\begin{figure*}[!t]
  \centering
  \begin{tikzpicture}[
  kairo/base,
  x=1cm,
  y=1cm,
  grid line/.style={draw=KairoLine, line width=0.45pt},
  value line/.style={line width=0.91pt, line cap=round}
]
  \node[anchor=east, inner sep=0pt, text=KairoMuted] at (0.08,2.32)
    {Score};
  \foreach \level/\tickLabel in {0.38/$0$,1.10/$\alpha$,1.46/$\tfrac32\alpha$} {
    \node[anchor=east, inner sep=0pt, text=KairoMuted] at (0.08,\level) {\tickLabel};
  }

  \foreach \fromX/\panelTitle/\ink in {
    0.24/{Previous window}/KairoBlue,
    4.65/{Recent mean}/KairoTeal,
    9.06/{Whole-log mean}/KairoMuted%
  } {
    \node[font=\footnotesize\bfseries, inner sep=0pt, text=\ink]
      at ({\fromX+2.02},2.32) {\panelTitle};
    \fill[KairoPaper] (\fromX,0.38) rectangle ++(4.04,1.63);
    \foreach \level in {1.10,1.46}
      \draw[grid line] (\fromX,\level) -- ++(4.04,0);
    \draw[KairoMuted, line width=0.5pt] (\fromX,0.38) -- ++(4.04,0);
    \foreach \offset/\windowName in {0.60/A,2.02/B,3.44/C} {
      \draw[KairoMuted, line width=0.45pt]
        ({\fromX+\offset},0.38) -- ++(0,-0.06);
      \node[anchor=north, text=KairoMuted, inner sep=0pt]
        at ({\fromX+\offset},0.25) {$b(V_{\mathrm{\windowName}})$};
    }
  }

  \draw[value line, KairoBlue] (2.26,1.10) -- (3.68,1.10);
  \foreach \x in {2.26,3.68}
    \fill[KairoBlue] (\x,1.10) circle[radius=2.0pt];
  \fill[KairoTeal] (8.09,1.46) circle[radius=2.0pt];
  \draw[value line, KairoMuted] (9.66,1.10) -- (11.08,0.38) -- (12.50,1.10);
  \foreach \x/\y in {9.66/1.10,11.08/0.38,12.50/1.10}
    \fill[KairoMuted] (\x,\y) circle[radius=2.0pt];
\end{tikzpicture}

  \caption{Illustrative total-variation scores. The scale uses the middle window's share $\alpha$ of state $\mathsf S$. Missing points indicate that the reference needs more history. The whole-log reference gives the middle window a score of zero.}
  \label{fig:approach-drift-scores}
\end{figure*}


A high score indicates an interval to inspect. It can reflect a change in process behavior or ordinary variation, particularly when a window has low support or covers only part of the monitored period. Different perspectives can respond to the same change. \Cref{sec:assessment-config} defines the benchmark's alarm rule for turning these scores into estimated change times.
